\section{Related Work}

Much of what this theory needs has already been proven by others. The work mostly \textbf{assembles} established results onto one specific substrate, the $\{3,4,3,4\}$ vibe mesh, and then asks one further question those programs leave aside. What is the felt interior of the thing, and which way does it move.

The aim of this section is honest credit. Each family below has settled some piece of the picture, and where it has, the theory leans on it rather than reinventing it. The section closes by stating, narrowly, what is genuinely new here.

A short orientation first, so the comparisons read cleanly. The universe of this theory is a vibe mesh, a flow of eight atoms. The mesh is the $\{3,4,3,4\}$ hyperbolic honeycomb. A dock is one cell, a here-now. A dock holds exactly $24$ sites, the $24$ directions out of it, and each site carries one tone, a value in $\{-1, 0, +1\}$ read as pain, peace, pleasure. The knit is the law of how the state changes, collide then stream. The wake is the growing edge where new docks are born. The beat is the discrete step. The lean is the value order $-1 < 0 < +1$, the single deliberate posit, and the arrow of time and the attraction between matter are emergent rather than base. With that vocabulary fixed, the families follow.

\subsection{Discrete spacetime}

Four programs take spacetime itself to be discrete and relational. The theory shares that starting point and differs on how each one secures Lorentz invariance, dynamics, and the status of the graph.

\textbf{Causal sets} (Bombelli, Lee, Meyer, Sorkin, with later work by Rideout, Sorkin, and Surya) model spacetime as a discrete relational order, with distance recovered by counting steps along causal chains. They buy Lorentz invariance with \textbf{randomness}, a Poisson sprinkling that has no preferred frame because it has no regular grid. The vibe mesh buys the same invariance with \textbf{curvature}. A flat lattice has a rest frame, but negative curvature scrambles directions so thoroughly that no global frame survives, and this is done while keeping the knit fully deterministic.

\textbf{Causal dynamical triangulations} (Ambjorn, Jurkiewicz, Loll) let geometry and even dimension emerge as outputs of a discrete dynamics, computed through a \textbf{path integral}, an ensemble average over triangulated histories. The vibe mesh fixes one rigid geometry and never sums over alternatives. The dimension is a property of the chosen honeycomb, not an expectation value.

\textbf{Loop quantum gravity and spin networks} (Rovelli, Smolin) treat space as a graph of relations with discrete capacity, where the graph itself is the fluctuating quantum state. In the vibe mesh the graph is fixed and the only dynamical variable is the ternary tone carried on the sites of each dock. The structure does not fluctuate, the values on it do.

\textbf{The Wolfram Physics Project} (Wolfram) is the closest neighbour in spirit. The universe is a graph plus a local rule, and time is rule application. There the hypergraph grows and rewrites itself, and a vast space of candidate rules is searched for one that fits. Here the graph is \textbf{fixed}, argued for rather than searched, and only the tone varies. Wolfram bets the structure is found by search. This theory bets it is forced by a few demands.

\begin{principle}[Determinism kept by curvature]
Where the discrete-spacetime programs reach Lorentz invariance through randomness or through summing over geometries, the vibe mesh reaches it through negative curvature alone, leaving the knit deterministic and the geometry single and fixed.
\end{principle}

\subsection{Deterministic quantum substrates}

A second family asks whether the quantum formalism rests on a deterministic classical layer. This is the family the theory leans on most heavily, since the knit is meant to \emph{be} such a layer.

\textbf{The cellular-automaton interpretation of quantum mechanics} ('t Hooft) argues that beneath the quantum formalism sits a deterministic, reversible, classical automaton. The theory inherits this conviction rather than originating it, and commits to one concrete substrate and one concrete knit.

\textbf{Quantum cellular automata and quantum walks to the Dirac equation} (D'Ariano and Perinotti, Bisio with D'Ariano and Tosini, Arrighi and collaborators, Farrelly and Short) supply the construction the theory most directly borrows. A unitary nearest-neighbour spin exchange spreads ballistically, and its continuum limit is the Dirac equation, with the local site angle setting the particle mass. This route is not discovered here. It belongs to these authors. The additions are the hyperbolic substrate, the classical-stochastic shadow that sits beneath the unitary layer, the forced geometry, and the observation that the knit \textbf{is} exactly this construction rather than an analogue of it.

\textbf{Quantum walks on Cayley graphs} (Lopez Acevedo and others) study walks on graphs generated by a Coxeter reflection group, which is to say Cayley-like graphs. The vibe mesh is generated this way, by reflections of a $\{3,4,3,4\}$ cell. The walk theory is borrowed. The hyperbolic Cayley arena, in place of the usual flat or finite ones, is the contribution.

\begin{result}[The knit as a known walk]
The collide-then-stream knit is, in its unitary reading, the spin-exchange quantum walk whose continuum limit is the Dirac equation. The theory does not derive this limit anew. It identifies the knit with the established construction and places it on a hyperbolic substrate.
\end{result}

\subsection{Holography and tensor networks}

\textbf{AdS/CFT and entanglement geometry} (Maldacena, Ryu and Takayanagi, Van Raamsdonk, Swingle's MERA construction, the HaPPY holographic code, Jahn and Eisert) tie bulk geometry to boundary entanglement, often as a static correspondence. These are largely fixed pictures, a code or a network laid down once. The vibe mesh lets a knit \textbf{run} and asks whether the area law, the bulk shortcut, and a working error-correcting code appear as \textbf{measured} properties of the dynamics rather than as built-in features. The result is a dynamical cousin of the static holographic codes, with the hyperbolic mesh supplying the bulk and the wake supplying the boundary in time.

\subsection{Hyperbolic lattices and automata}

A third family studies hyperbolic lattices directly, in theory and in the laboratory.

\textbf{Hyperbolic lattices in the lab} (Kollar, Fitzpatrick, and Houck in circuit quantum electrodynamics, Maciejko and Rayan in hyperbolic band theory) show that the geometry the theory bets on is one laboratories can build and probe. They study one phenomenon against a hyperbolic backdrop. The theory proposes the lattice as the substrate of everything.

\textbf{Hyperbolic cellular automata} (Margenstern) establish universality across hyperbolic spaces, through the railway model, and develop Fibonacci-style addressing for tilings such as $\{7,3\}$ and $\{5,3,4\}$. The theory needs universality on its \textbf{own} knit, which the functional-completeness and Minsky-machine results confirm, and treats $\{3,4,3,4\}$ addressing as an open extension of Margenstern's program.

\textbf{Hyperbolic graphs and geometric group theory} (Cannon, Hamann) prove that negatively curved graphs are expanders, tree-like at large scale, with exponential ball growth. The theory turns these graph-theoretic facts into physical consequences. The expansion gives a speed limit, the large-scale tree structure gives an emergent metric, the fast mixing gives quick integration, and the boundary at infinity gives a holographic readout.

\begin{proposition}[Geometry into physics]
The expander, tree-like, exponentially growing structure of negatively curved graphs, established by geometric group theory, becomes in the vibe mesh a speed limit, an emergent metric, fast integration, and a holographic readout.
\end{proposition}

\subsection{Associative and content-addressable computing}

\textbf{Associative computing} (Potter) treats memory as content-addressable and computes by broadcasting a query to every cell at once, each cell answering locally with no central scan. This is exactly the model the theory's associative-memory engine realizes on the hyperbolic bulk, where recall is lookup by content and the parallel match is a single broadcast. The paradigm is borrowed. The contribution is that the negative curvature of the substrate makes that broadcast logarithmic in the size of the memory, where a flat array would be polynomial. The exponential room of the hyperbolic mesh reaches every stored pattern in a number of beats that grows only with the logarithm of how many there are.

\subsection{Foundations and mind}

\textbf{Informational reconstructions of quantum mechanics} (Chiribella, D'Ariano, Perinotti) argue that quantum structure should be \textbf{derived} from information-theoretic axioms, not postulated, and ideally without committing to any substrate. The theory takes the opposite road. It commits to a substrate and lets the structure fall out. The two approaches are complementary, the axiomatic why and the constructive how.

\textbf{Consciousness} (Chalmers' hard problem, Tononi's Integrated Information Theory, Wheeler's slogan that physics comes from information) frames the question the theory answers by monism. The theory takes experience as the \textbf{substance} of the substrate, the vibe itself, rather than a quantity computed from structure after the fact. This sits close to idealism and to Wheeler's informational starting point, and it uses integration, in the spirit of Integrated Information Theory, to address the combination problem that panpsychism faces, namely how many small experiences add up to one.

\subsection{Emergent and entropic gravity}

\textbf{Gravity as thermodynamics} (Jacobson, Padmanabhan, Verlinde) argues that gravity is not fundamental. Jacobson derives the Einstein equation as an equation of \textbf{state}, from heat and area at a local horizon. Verlinde recovers Newton's law, with readings that reach the dark-universe and MOND scales, from entropy and information. This is the natural home for the theory's \textbf{result} that the $1/r$ gravitational law is not fundamental but emerges from the substrate's geometry and the conservation built into the knit. The theory derives the $1/r$ directly on the mesh. These authors argue, independently, that gravity-as-thermodynamics is exactly where such a derivation belongs.

\begin{result}[Emergent inverse-square law]
The $1/r$ gravitational law is derived on the vibe mesh from geometry and the conservation of the knit, not assumed. The entropic-gravity programs supply the prior argument that gravity is the kind of thing that should emerge this way.
\end{result}

\subsection{Positive geometry, a forward-looking ally}

\textbf{The amplituhedron and positive geometries} (Arkani-Hamed and Trnka, Arkani-Hamed with Bai and Lam) hold that physical structure is the \textbf{output} of a primitive shape together with a \textbf{positivity} condition. The amplituhedron's recursive-boundary structure mirrors the theory's holographic code and its coarse-graining tower, where each level is the boundary of the one below. The resonance is philosophical, geometry plus positivity yields physics, and it is flagged here as an alliance to pursue rather than a result already in hand.

\subsection{Mind, agency, and the strongest adversary}

The last family bears on the self, on agency, and on the sharpest objection to a deterministic substrate.

\textbf{Active inference} (Friston) supplies the working model of the self. In the theory the self \textbf{is} active inference. Its forward model is perception, its action is will, and the lean is its preference, the value order that tells it which states to seek.

\textbf{Causal emergence} (Hoel) proves that a \textbf{macro} level can carry more causal power than the micro level beneath it. This is the rigorous backbone for treating the self as a genuine author of its actions, and for the macro-compatibilist account of freedom the theory adopts.

\textbf{The free-will theorems} (Conway and Kochen) are the sharpest \textbf{adversary} to a deterministic, no-randomness substrate. The theory names the theorem rather than dodging it. It concedes the micro-deterministic door the theorem points to, and claims freedom only at the \textbf{macro} level, which is the honest move the theorem forces.

\textbf{Origin-of-life physics} (England's dissipative adaptation, Walker and Davies, assembly theory) offers candidate physics for why the lean favours reproducing, organized patterns, why life is a thing the substrate tends to make rather than an accident.

\textbf{Constructor theory} (Deutsch and Marletto) gives the natural framing for the level-relative question of which tone transformations are possible and which are impossible, a question the theory inherits when it asks what a self at one level can do to the tones at another.

\subsection{What is genuinely ours}

Stated narrowly, the following are the theory's own contributions.

\begin{enumerate}
\item A substrate argued to be \textbf{selected} rather than random, rewritten, ensemble-averaged, or assumed. This is a reduction of arbitrariness, not a derivation from nothing.
\item The ternary tone together with the \textbf{lean} as a value direction, the single deliberate posit of the whole framework.
\item One substrate that aims at physics \textbf{and} life \textbf{and} mind, rather than at any one of them alone.
\item Experiential \textbf{monism}, that the vibe is the substance and not a quantity computed from it, the part furthest outside the existing literature.
\item A falsifiable single-substrate program that states its negatives openly, the phenomena that did not emerge as well as those that did.
\end{enumerate}

The selection argument that grounds the substrate, the dimension ladder together with the $D_4$ spinor requirement, is a strong basis. The $\{3,4,3,4\}$ honeycomb alone is crystallographic and hyperbolic, hands over a three-dimensional flat cusp, and carries spinor sites in its $24$ directions. The theory leans on exactly that combination of facts.

\begin{principle}[The narrow claim]
The literature already proves most of the parts. What is new is the selection of one substrate, the lean as the one value posit, the single mesh aimed at physics and life and mind together, and the monism that takes the vibe as substance rather than as a derived measure.
\end{principle}

\subsection*{See also}

The companion sections develop what follows from the program. Implications takes up what holds if the theory is right. Scope and Coverage marks where the edges are and what is left undone. The Conclusion names the direction worth exploring next.
